%%%%%%%%%%%%%%%%%%%%%%%%%%%%%%%%%%%%%%%%%%%%%%%%%%%%%%%%%%%%%%%%%%%%%%%%%%%%%%%
%                         CAPITULO 4
%                  Unidad 4: Servicios de Red
%%%%%%%%%%%%%%%%%%%%%%%%%%%%%%%%%%%%%%%%%%%%%%%%%%%%%%%%%%%%%%%%%%%%%%%%%%%%%%%

\chapter{Servicios de Red}

Hasta ahora hemos aprendido como los dispositivos se conectan fisicamente
mediante cables y ondas de radio, como se organizan en redes, como se
identifican mediante direcciones IP y como los routers encaminan los paquetes.
Pero una red no es solo cableado y protocolos de bajo nivel: necesita
\textbf{servicios} que permitan a los usuarios y aplicaciones compartir recursos,
comunicarse y acceder a internet.

En este capitulo estudiaremos los principales servicios de red, con especial
atencion a dos de los mas fundamentales: el \textbf{DNS} (sistema de nombres de
dominio) y el \textbf{DHCP} (protocolo de configuracion dinamica de host), sin
los cuales una red moderna no podria funcionar.


\section{Definicion de servicios de red}

\begin{definicion}{Servicios de red}
Los \textbf{servicios de red} son protocolos de software ejecutados en servidores,
que permiten a los dispositivos conectados compartir recursos, comunicarse,
acceder a internet y gestionar la infraestructura tecnologica.
\end{definicion}

Los servicios de red operan bajo un modelo \textbf{cliente-servidor}: el servidor
ofrece un servicio y el cliente lo consume. Son esenciales tanto en redes
locales (LAN) para compartir impresoras o archivos, como en redes globales
para el acceso a internet.


\section{Clasificacion de los servicios de red}

Los servicios de red se pueden clasificar en cuatro categorias principales:

\subsection{Servicios de infraestructura}

Son los servicios fundamentales para la conectividad y el funcionamiento de la
red:

\begin{itemize}
    \item \textbf{DNS (Domain Name System)}: traduce nombres de dominio
    (direcciones web como \texttt{www.upv.es}) a direcciones IP numericas
    (como \texttt{158.42.4.23}). Es el ``agenda telefonica'' de Internet.
    \item \textbf{DHCP (Dynamic Host Configuration Protocol)}: asigna
    automaticamente direcciones IP y otros parametros de configuracion (mascara,
    gateway, DNS) a los dispositivos que se conectan a la red.
\end{itemize}

\subsection{Servicios de internet y comunicacion}

Permiten la comunicacion entre personas y el acceso a la informacion:

\begin{itemize}
    \item \textbf{WWW (World Wide Web)}: usa los protocolos HTTP/HTTPS para
    navegar por paginas web. Es el servicio mas visible de Internet.
    \item \textbf{Correo electronico}: usa SMTP (envio), POP3 e IMAP (recepcion)
    para el intercambio de mensajes.
    \item \textbf{FTP (File Transfer Protocol)}: permite transferir ficheros entre
    ordenadores de forma rapida y fiable.
\end{itemize}

\subsection{Servicios de acceso y gestion}

Facilitan la administracion remota y la organizacion de recursos:

\begin{itemize}
    \item \textbf{Acceso remoto}: SSH (Secure Shell) y Telnet permiten gestionar
    equipos a distancia. SSH es el estandar actual porque cifra la comunicacion.
    \item \textbf{Servicios de directorio}: gestionan usuarios, permisos y recursos
    en una red. Ejemplo: Active Directory de Microsoft.
\end{itemize}

\subsection{Servicios de aplicacion}

Son los servicios que usan directamente los usuarios finales:

\begin{itemize}
    \item \textbf{Aplicaciones en la nube} (SaaS): software ejecutado en servidores
    remotos (Google Workspace, Microsoft 365).
    \item \textbf{Mensajeria}: WhatsApp, Slack, Teams.
    \item \textbf{Bases de datos}: servidores que almacenan y gestionan datos.
    \item \textbf{Streaming}: Netflix, Spotify, YouTube.
\end{itemize}


\section{DNS (Domain Name System)}

\subsection{El problema: los humanos usan nombres, las maquinas usan numeros}

Los seres humanos recordamos nombres mucho mejor que numeros. Sin embargo,
los dispositivos de red se identifican mediante direcciones IP (numeros). El
\textbf{DNS} resuelve esta disparidad actuando como un traductor entre nombres y
numeros.

\begin{definicion}{DNS (Domain Name System)}
El \textbf{DNS} (Domain Name System) es un sistema distribuido y jerarquico que
traduce nombres de dominio legibles por humanos (como \texttt{www.upv.es}) a
direcciones IP numericas (como \texttt{158.42.4.23}) y viceversa. Esta definido en
los RFC 1034 y 1035.
\end{definicion}

\subsection{El nombre en una red: FQDN}

Un nombre de dominio completo se llama \textbf{FQDN} (Full Qualified Domain Name)
y tiene una estructura jerarquica:

\begin{verbatim}
www.upv.es
|    |   |
|    |   +-- Dominio de nivel superior (TLD): pais o categoria
|    +------ Organizacion a la que pertenece
+----------- Nombre del equipo o rol que desempena
\end{verbatim}

Por ejemplo, en \texttt{www.upv.es}:
\begin{itemize}
    \item \texttt{www}: nombre del equipo o servicio (world wide web)
    \item \texttt{upv}: organizacion (Universitat Politecnica de Valencia)
    \item \texttt{es}: codigo de pais (Espana)
\end{itemize}

\subsection{Estructura jerarquica del DNS}

El DNS no es una base de datos centralizada. Es \textbf{distribuido} en una
estructura de arbol invertido:

\begin{enumerate}
    \item \textbf{Root Servers} (servidores raiz): son los 13 servidores raiz
    etiquetados de la A a la M. La mayoria estan en Estados Unidos. Son el
    punto de entrada de cualquier consulta DNS.
    \item \textbf{TLD Servers} (Top Level Domain): gestionan los dominios de
    primer nivel (.com, .es, .org, .edu, .net).
    \item \textbf{Authoritative DNS Servers}: servidores autoritativos de cada
    dominio concreto. Contienen los registros reales (A, MX, CNAME...) de
    los equipos de esa organizacion.
\end{enumerate}

Esta descentralizacion evita un punto unico de fallo y distribuye el trafico.

\subsection{Tipos de consultas DNS}

Existen dos formas de resolver un nombre:

\begin{itemize}
    \item \textbf{Consulta recursiva}: el cliente pide al servidor DNS que resuelva
    el nombre completo. Si el servidor no lo sabe, el propio servidor consulta
    a otros servidores en nombre del cliente hasta obtener la respuesta.
    \item \textbf{Consulta iterativa}: el servidor responde con la mejor
    informacion que tiene. Si no conoce la respuesta, indica al cliente a que
    otro servidor debe preguntar.
\end{itemize}

\textbf{El caching (cache) es fundamental}: los servidores y los clientes
almacenan temporalmente las respuestas para acelerar futuras consultas y
reducir la carga en la red.

\subsection{Tipos de registros DNS}

Un servidor DNS autoritativo almacena registros de distintos tipos:

\begin{table}[htbp]
\centering
\small
\caption{Tipos de registros DNS mas comunes}
\begin{tabularx}{\textwidth}{@{}l>{\raggedright\arraybackslash}X@{}}
\toprule
\textbf{Registro} & \textbf{Funcion} \\
\midrule
A (Address) & Asocia un nombre de dominio a una direccion IPv4 \\
AAAA & Asocia un nombre a una direccion IPv6 \\
CNAME (Canonical Name) & Alias: asocia un nombre alternativo a un nombre canonico \\
MX (Mail Exchange) & Indica el servidor de correo del dominio \\
TXT & Texto arbitrario, usado para verificaciones y politicas (SPF, DKIM) \\
NS (Name Server) & Indica que servidores DNS son autoritativos para el dominio \\
PTR (Pointer) & Resolucion inversa: de IP a nombre de dominio \\
\bottomrule
\end{tabularx}
\end{table}

\subsection{DNS internos vs. DNS externos (autoritativos)}

\begin{itemize}
    \item \textbf{DNS interno}: resuelve nombres dentro de la red local de una
    empresa. Por ejemplo, \texttt{servidor-interno.empresa.local} puede apuntar a
    \texttt{192.168.1.10}. Solo es visible desde dentro de la red.
    \item \textbf{DNS externo / autoritativo}: resuelve nombres publicos en
    Internet. Por ejemplo, \texttt{www.empresa.com} apunta a la IP publica del
    servidor web. Es visible desde cualquier parte de Internet.
\end{itemize}

Muchas empresas tienen ambos: un DNS interno para sus equipos y un DNS
externo gestionado por su proveedor o registrador de dominio.

\subsection{Consulta manual con nslookup}

Para diagnosticar problemas DNS, se puede usar el comando \texttt{nslookup}:

\begin{lstlisting}[language=bash,caption={Consulta DNS con nslookup},label={lst:nslookup}]
C:\> nslookup www.upv.es
Servidor:  dns.google
Address:  8.8.8.8

Respuesta no autoritativa:
Nombre:  www.upv.es
Address:  158.42.4.23
\end{lstlisting}


\section{DHCP (Dynamic Host Configuration Protocol)}

\subsection{Definicion}

\begin{definicion}{DHCP (Dynamic Host Configuration Protocol)}
El \textbf{DHCP} es un protocolo de red que permite a los dispositivos obtener
automaticamente una direccion IP y otros parametros de configuracion (mascara de
subred, gateway, servidores DNS) de un servidor, sin intervencion manual.
\end{definicion}

Sin DHCP, cada dispositivo de una red tendria que configurarse manualmente con
una IP, mascara, gateway y DNS. En una red con cientos de dispositivos, esto seria
impracticable y propenso a errores (IPs duplicadas, mascaras incorrectas...).

\subsection{Beneficios clave de DHCP}

\begin{itemize}
    \item \textbf{Administracion centralizada}: un unico servidor gestiona toda la
    configuracion IP de la red.
    \item \textbf{Consistencia}: todos los dispositivos reciben la configuracion
    correcta y homogenea.
    \item \textbf{Flexibilidad}: cambiar la red (por ejemplo, el gateway) solo requiere
    modificar el servidor, no cada dispositivo.
    \item \textbf{Escalabilidad}: anadir nuevos dispositivos es tan simple como
    conectarlos; reciben IP automaticamente.
\end{itemize}

\subsection{Como funciona: el proceso DORA}

Cuando un dispositivo se conecta a una red, el proceso DHCP sigue cuatro pasos:

\begin{enumerate}
    \item \textbf{DHCPDISCOVER}: el cliente envia un mensaje de difusion (broadcast)
    preguntando: ``Hay algun servidor DHCP? Necesito una IP''.
    \item \textbf{DHCPOFFER}: el servidor responde ofreciendo una direccion IP
    disponible, junto con la mascara, el gateway, los DNS y el tiempo de
    concesion (lease time).
    \item \textbf{DHCPREQUEST}: el cliente acepta la oferta y solicita formalmente
    esa IP.
    \item \textbf{DHCPACK}: el servidor confirma la concesion y el cliente puede
    empezar a usar la IP.
\end{enumerate}

Este proceso se conoce como \textbf{DORA} (Discover, Offer, Request, Acknowledge).

\begin{figure}[H]
\centering
\fbox{\parbox{0.8\textwidth}{\centering\vspace{1.5cm}\textit{TODO: Anadir diagrama del proceso DORA (Discover, Offer, Request, Acknowledge).}\vspace{1.5cm}}}
\caption{Proceso DORA de DHCP}
\label{fig:dhcp-dora}
\end{figure}

\subsection{Mensajes DHCP adicionales}

Ademas de los cuatro principales, existen otros mensajes:

\begin{itemize}
    \item \textbf{DHCPDECLINE}: el cliente rechaza la IP ofrecida porque detecta
    que ya esta en uso (por ejemplo, mediante ARP).
    \item \textbf{DHCPNAK}: el servidor deniega una concesion (por ejemplo, si la IP
    ya no esta disponible o el cliente ha cambiado de red).
    \item \textbf{DHCPRELEASE}: el cliente libera voluntariamente la IP antes de
    desconectarse (util para dispositivos moviles).
    \item \textbf{DHCPINFORM}: el cliente ya tiene IP (por ejemplo, configurada
    manualmente) pero solicita otros parametros como DNS o gateway.
\end{itemize}

\subsection{Renovacion de la concesion}

Las IPs asignadas por DHCP no son permanentes: se conceden por un tiempo
limitado (\textbf{lease time}). Cuando la mitad del tiempo ha transcurrido, el
cliente intenta renovar la concesion enviando un nuevo DHCPREQUEST. Si el
servidor responde con DHCPACK, la concesion se extiende. Si no, el cliente
puede seguir usando la IP hasta que expire, momento en el que debe iniciar un
nuevo proceso DORA.

\subsection{Ambitos (scopes) y exclusiones}

Un servidor DHCP organiza las direcciones en \textbf{ambitos} (scopes):

\begin{itemize}
    \item \textbf{Ambito}: un rango continuo de direcciones IP disponibles para
    asignar. Por ejemplo, de 192.168.1.100 a 192.168.1.200.
    \item \textbf{Exclusiones}: direcciones dentro del ambito que no se asignan
    automaticamente. Por ejemplo, la IP del router (192.168.1.1), del servidor
    de impresion (192.168.1.10), etc.
    \item \textbf{Reservas}: direcciones asignadas permanentemente a un dispositivo
    especifico (identificado por su direccion MAC). Util para servidores,
    impresoras o PLCs que necesitan una IP fija.
\end{itemize}

\subsection{Agentes de retransmision (DHCP Relay)}

Si el servidor DHCP esta en una red diferente a la de los clientes, los mensajes
DHCP (que son broadcasts) no llegan a traves de los routers. La solucion son los
\textbf{agentes de retransmision} (DHCP relay agents): dispositivos que reciben los
mensajes DHCP de broadcast en una red y los reenvian como unicast al servidor
DHCP en otra red.

\subsection{BOOTP: el predecesor de DHCP}

\textbf{BOOTP} (Bootstrap Protocol) fue el predecesor de DHCP. A diferencia de
DHCP, BOOTP asignaba direcciones IP de forma estatica (manualmente
configuradas en el servidor) y no admitia concesiones temporales. DHCP es
retrocompatible con BOOTP y puede atender a clientes BOOTP.


\section{Otros servicios de red relevantes}

\subsection{HTTP/HTTPS (World Wide Web)}

El protocolo \textbf{HTTP} (HyperText Transfer Protocol) es el fundamento de la Web.
Permite a los navegadores (clientes) solicitar paginas web a los servidores.
\textbf{HTTPS} es la version segura, que cifra la comunicacion mediante TLS/SSL.

\subsection{SSH (Secure Shell)}

\textbf{SSH} permite acceder y gestionar equipos remotos mediante una conexion
segura y cifrada. Es el reemplazo seguro de Telnet (que enviaba datos en texto
plano). Se usa ampliamente para administrar servidores, routers, switches y
dispositivos IoT.

\subsection{Correo electronico: SMTP, POP3, IMAP}

\begin{itemize}
    \item \textbf{SMTP} (Simple Mail Transfer Protocol): envia correos desde el
    cliente al servidor, y entre servidores.
    \item \textbf{POP3} (Post Office Protocol): descarga los correos del servidor al
    cliente (los elimina del servidor).
    \item \textbf{IMAP} (Internet Message Access Protocol): accede a los correos en
    el servidor sin descargarlos, permitiendo sincronizacion entre multiples
    dispositivos.
\end{itemize}

\subsection{FTP/SFTP}

\textbf{FTP} (File Transfer Protocol) permite transferir ficheros entre ordenadores.
Aunque es rapido, no cifra la comunicacion. \textbf{SFTP} (SSH File Transfer Protocol)
utiliza SSH para transferir ficheros de forma segura.


\section{Resumen de la Unidad 4}

\begin{itemize}
    \item Los \textbf{servicios de red} son protocolos de software que permiten a los
    dispositivos compartir recursos, comunicarse y gestionar la infraestructura.
    \item Los principales tipos de servicios son: infraestructura (DNS, DHCP),
    internet y comunicacion (WWW, correo, FTP), acceso y gestion (SSH, Active
    Directory), y aplicacion (nube, bases de datos, streaming).
    \item El \textbf{DNS} traduce nombres de dominio a direcciones IP. Es un sistema
    jerarquico y distribuido con servidores raiz, TLD y autoritativos. Utiliza
    registros como A, CNAME, MX, TXT, NS.
    \item El \textbf{DHCP} asigna automaticamente direcciones IP y parametros de
    configuracion. El proceso sigue cuatro pasos: Discover, Offer, Request,
    Acknowledge (DORA). Gestiona ambitos, exclusiones, reservas y renovaciones.
    \item La \textbf{convergencia de servicios} (DNS + DHCP + Active Directory) permite
    una gestion centralizada y automatizada de las redes modernas.
\end{itemize}
